% Figure from MATH 590 notes, section 16. Compile with the notes preamble.
\begin{tikzpicture}[x=0.82cm, y=1cm]
  \foreach \n in {1,...,6} {
    \draw[figB, dashed, rounded corners=5pt, fill=figB!8] (\n-0.3,-0.35) rectangle (\n+0.3,1.35);
    \node[font=\scriptsize, below=1pt] at (\n,-0.35) {$\n$};
    \fill[black!45] (\n,0) circle (1.5pt) (\n,1) circle (1.5pt);
  }
  \node[font=\scriptsize] at (6.85,0.5) {$\cdots$};
  \node[font=\scriptsize, left] at (0.6,0) {$a$};
  \node[font=\scriptsize, left] at (0.6,1) {$b$};
  \node[font=\scriptsize, figB] at (3.5,1.72) {open sets $\{n\}\times Y$: no finite subfamily covers $X$};
  \node[font=\scriptsize] at (6.85,-0.62) {$\mathbb{Z}_+$};
  % an infinite set S (red) and one of its limit points
  \foreach \p in {(1,0),(2,1),(3,0),(4,0),(5,1),(6,0)} \fill[figR] \p circle (2pt);
  \draw[figR, thick] (3,1) circle (3.6pt);
  \node[font=\scriptsize, figR, align=left, right] at (7.35,0.9) {$3\times b$: limit point\\ of $S\ni 3\times a$};
  \node[font=\scriptsize, figR, right] at (7.35,0.0) {points of $S$};
\end{tikzpicture}
